\section{Discussion}
\label{sec:discussion}

\subsection{Key Findings}

\textbf{Finding 1:} Behavioral proxy signals are detectable in returning-user cohorts, but in the wrong direction for BAA disengagement. The positive prompt length trend ($\tau = +0.744$) establishes that behavioral signals are computationally accessible --- but contradicts the BAA prediction.

\textbf{Finding 2:} Two of three proxies are unmeasurable under current conditions. Vote entropy is blocked by dataset access; correction frequency has too sparse a signal. Only one proxy is evaluable, and it shows the opposite trend.

\textbf{Finding 3:} The returning-user cohort design introduces a selection bias that prevents causal attribution. The $\geq$3 monthly bin filter is necessary for signal detection but systematically retains high-engagement users.

\subsection{Why Is Prompt Length Increasing?}

Three competing explanations are consistent with the positive trend, all independent of AI quality improvement. This paper cannot discriminate among them with the available data; the specific test required to do so is a difference-in-differences comparison between returning and non-returning users (described in Section~\ref{sec:limitations}, L1).

\textbf{Explanation 1: Selection bias (most plausible).} High-engagement users who return $\geq$3 months are power users who compose longer prompts by baseline behavioral profile. The positive trend may reflect increasing dominance of heavy users over time. \emph{Discriminating test:} compare prompt length trajectory for users in their first appearance month vs.\ users in their third+ appearance month.

\textbf{Explanation 2: User expertise gain (also plausible).} Users learn to compose more effective prompts over time --- longer, more context-rich. This expertise gain predicts increasing prompt length independent of AI quality. \emph{Discriminating test:} within-user prompt length trajectory (requires individual-level longitudinal data not available in this study).

\textbf{Explanation 3: Platform adoption effect (moderate).} WildChat's user base shifted toward professional/developer use, involving longer technical prompts. \emph{Discriminating test:} topic-tag stratification of the prompt length trend.

The BAA disengagement mechanism requires decreasing prompt length. Both selection bias and expertise gain offer more parsimonious explanations. The paper cannot rule out any of these; future work with the comparison group described in L1 can.

\subsection{Limitations}
\label{sec:limitations}

\textbf{L1: Returning-user selection bias (severity: HIGH).} The analysis cohort ($n = 27{,}902$) is self-selected high-engagement. The proposed control experiment: non-returning user comparison in difference-in-differences design.

\textbf{L2: LMSYS primary dataset access (severity: HIGH for P2).} Vote entropy --- arguably the most direct BAA operationalization --- is entirely blocked. Remediation: request LMSYS research access.

\textbf{L3: Correction frequency proxy inadequacy (severity: MEDIUM).} Regex captures explicit corrections ($\approx$0.05\% of turns). Implicit corrections require LLM-based annotation or session-level proxies.

\textbf{L4: No causal chain verification (severity: HIGH).} The BAA mechanism posits three steps: (1) AI quality improves, (2) users perceive less need to probe carefully, (3) behavioral engagement declines. None of these steps is verified here. Most critically, WildChat contains no per-interaction measure of AI model quality or version; we cannot link observed prompt length trends to specific periods of AI improvement. The empirical result ($\tau = +0.744$) is informative about the direction of behavioral change in the observed cohort --- but it cannot confirm or refute the BAA causal mechanism. Future work must include an AI quality covariate (e.g., per-bin model version or quality rating) to test the causal chain.

\textbf{L5: IP-hash cohort noise (severity: MEDIUM).} Multiple users may share an IP (NAT, shared networks); a single user may appear as multiple IPs (VPN, mobile). The cohort ($n = 27{,}902$) is defined by IP-hash continuity. The magnitude of this noise is not quantified.

\textbf{L6: WildChat platform scope (severity: LOW--MEDIUM).} WildChat logs are ChatGPT (OpenAI API) interactions only. Results may not generalize to other AI platforms (Claude, Gemini, open-source model deployments).

\textbf{L7: Model transition confound (severity: MEDIUM).} GPT-4 API launched in March 2023 --- immediately before our observation window (April 2023). Users switching to GPT-4, which supports substantially longer context windows and better instruction following than GPT-3.5-turbo, may have begun composing longer prompts for technical rather than behavioral reasons: GPT-4 can handle longer, more detailed inputs effectively, incentivizing longer prompt composition regardless of user engagement level. This confound partially overlaps with L4 (no AI quality covariate) but is more specific: the GPT-3.5 $\rightarrow$ GPT-4 transition is a discrete event whose timing aligns with the start of our observation window. Disentangling model capability effects from user behavioral adaptation requires holding the AI model constant across the observation period, which is not possible in WildChat-1M.

\subsection{Broader Impact}

This work develops methods for longitudinal behavioral analysis in large-scale AI interaction logs. The analysis pipeline (returning-user cohort construction, Hamed-Rao Mann-Kendall, bootstrap CI) is applicable to any dataset with user-level pseudonymization and temporal coverage. Our analysis operates exclusively at the cohort aggregate level; the IP-hash approach does not enable re-identification.
